%% ma: L10-B12-0014
%% lop: 10
%% chuong: 5
%% bai: 12
%% dang: 10.12.4
%% loai: TL
%% muc: TH
%% dap_an: "-"
%% nguon: "chuyen-de-cac-so-dac-trung-cua-mau-so-lieu-khong-ghep-nhom-toan-10-chuong-trinh-moi.pdf (GV. Phan Nhật Linh) (P01), Bài 01 Bài tập tự luận 14"
%% kien_thuc: "Lớp 10 Bài 12 (số gần đúng, sai số tuyệt đối, độ chính xác, sai số tương đối, quy tròn số)"
%% trang_thai: chinh-thuc
%% ghi_chu: "Tự duyệt theo yêu cầu giáo viên 22/10/2026."
\begin{ex}
Các nhà khoa học đo được khoảng cách từ Trái Đất đến Sao Hỏa là $55$ triệu km $\pm\,43$ km, từ Trái Đất đến Sao Kim là $38$ triệu km $\pm\,31$ km, từ Trái Đất đến Mặt Trời là $150$ triệu km $\pm\,101$ km. Phép đo nào chính xác nhất?
\loigiai{Sao Hỏa: $\delta_1\le\dfrac{43}{55\cdot10^6}\approx7{,}82\cdot10^{-7}$.

Sao Kim: $\delta_2\le\dfrac{31}{38\cdot10^6}\approx8{,}16\cdot10^{-7}$.

Mặt Trời: $\delta_3\le\dfrac{101}{150\cdot10^6}\approx6{,}73\cdot10^{-7}$.

Vậy phép đo khoảng cách từ Trái Đất đến Mặt Trời chính xác nhất.}
\end{ex}
